\section*{Required author Section 5.1: diagonalization of the expectation}
Throughout Section5, $M'$ denotes the $\binom{[n]}{d/2}\times\binom{[n]}{d/2}$ minor of the non-exact moment matrix:
\[M'(I,J)=\sum_{T:|V(T)\cup I\cup J|\le\tau}(\omega/n)^{|V(T)\cup I\cup J|}\chi_T,\quad\forall I,J\in\binom{[n]}{d/2}.\tag{5.1}\]
\footnote{Strictly speaking PSDness of this minor is not sufficient since there is no homogeneity reduction in the non-exact case; it suffices to demonstrate the idea.} The goal is to diagonalize $M'$ approximately so that the difference matrix is negligible w.h.p. on plugging $G$.
\paragraph{Proposition 5.1.} $\mathbb E[M']=CC^\top$, where $C$ is the $\binom{[n]}{d/2}\times\binom{[n]}{\le d/2}$ matrix
\[C=(\zeta^\top)_{d/2,\le d/2}\operatorname{diag}(\sqrt{t(|A|)})_{A\in\binom{[n]}{\le d/2}},\tag{5.2}\]
and $t(r)=(1-O(d\omega/n))(\omega/n)^{d-r}$ for $r=0,\ldots,d/2$.
\paragraph{Definition 5.2 (See e.g. [9]).} Fix $n,k$. A Johnson scheme $J$ is a $\binom{[n]}k\times\binom{[n]}k$ matrix satisfying $J(I,J)=J(I',J')$ whenever $|I\cap J|=|I'\cap J'|$. For fixed $n,k$, Johnson schemes are symmetric matrices and form a commutative $\mathbb R$-algebra, hence are simultaneously diagonalizable. Fix $k=d/2$. One basis is $D_0,\ldots,D_{d/2}$ where
\[D_r(I,J)=\begin{cases}1,&|I\cap J|=r,\\0,&\text{o.w.}\end{cases}\tag{5.3}\]
Another basis is $J_0,\ldots,J_{d/2}$,
\[J_r(I,J)=\binom{|I\cap J|}r.\tag{5.4}\]
These are PSD since
\[J_r=\sum_{A\subseteq[n],|A|=r}u_Au_A^\top,\quad u_A\in\mathbb R^{\binom{[n]}k},\ u_A(B)=\mathbf1_{A\subseteq B}.\tag{5.5}\]
Clearly $J_{d/2}=\mathrm{Id}$.
\paragraph{Fact 5.3 (See e.g. (4.29) in [9]).} The Johnson schemes for $(n,d/2)$ have shared eigenspace-decomposition $\mathbb R^{\binom{[n]}{d/2}}=V_0\oplus\cdots\oplus V_{d/2}$, and $J_r=\bigoplus_{i=0}^{d/2}\lambda_r(i)\Pi_i$, where $\Pi_i$ is orthogonal projection to $V_i$ w.r.t. Euclidean inner product, and
\[\lambda_r(i)=\binom{d/2-i}{r-i}\binom{n-d/2-i}{d/2-r},\qquad0\le i\le d/2.\]
\paragraph{Lemma 5.4.} $\mathbb E[M']=\sum_{r=0}^{d/2}t(r)J_r$, with $t(r)=(1-O(d\omega/n))(\omega/n)^{d-r}$.
\begin{proof}By definition $\mathbb E[M']=\sum_{r=0}^{d/2}(\omega/n)^{d-r}D_r$. Each $D_r$ decomposes as
\[D_r=\sum_{r'=r}^{d/2}(-1)^{r'-r}\binom{r'}rJ_{r'}.\tag{5.6}\]
Indeed the RHS at $(I,J)$ equals
\begin{align*}&\sum_{r'=r}^{d/2}(-1)^{r'-r}\binom{r'}r\binom{|I\cap J|}{r'}\\&=\sum_{r'=r}^{|I\cap J|}(-1)^{r'-r}\binom{|I\cap J|}r\binom{|I\cap J|-r}{r'-r}\\&=\binom{|I\cap J|}r\mathbf1_{|I\cap J|=r}=\mathbf1_{|I\cap J|=r}.\end{align*}
So
\begin{align*}\mathbb E[M']&=\sum_{r=0}^{d/2}(\omega/n)^{d-r}\left(\sum_{r'=r}^{d/2}(-1)^{r'-r}\binom{r'}rJ_{r'}\right)\\&=\sum_{r'=0}^{d/2}J_{r'}\left(\sum_{r=0}^{r'}(\omega/n)^{d-r}(-1)^{r'-r}\binom{r'}r\right)\\&=\sum_{r'=0}^{d/2}J_{r'}(\omega/n)^{d-r'}(1-\omega/n)^{r'}.\tag{5.7}\end{align*}
This proves the lemma.\end{proof}
By Lemma5.4 and (5.5), letting $t(r')=(\omega/n)^{d-r'}(1-\omega/n)^{r'}$, we have
\begin{align*}\mathbb E[M']&=\sum_{A:|A|\le d/2}t(|A|)u_Au_A^\top\\&=(\zeta^\top)_{d/2,\le d/2}\operatorname{diag}(t(|A|))\zeta_{\le d/2,d/2}=CC^\top,\end{align*}
since $(\zeta^\top)_{d/2,\le d/2}$ has columns $\{u_A:|A|\le d/2\}$. This proves Proposition5.1.
\paragraph{Fact 5.7 ([11], required separator notation from Section5.2).} For ribbon $(I,J;T)$, all min-separators $\mathrm{mSep}_{I,J}(T)$ have a natural poset structure: $A_1\le A_2$ iff $A_1$ separates $(I,A_2;T)$, equivalently iff $A_2$ separates $(J,A_1;T)$. This poset is a lattice: each two min-separators have join and meet. In particular it has a unique minimum and maximum, denoted $S_l(I,J;T)$ and $S_r(I,J;T)$, the leftmost and rightmost min-separator. The rest of Section5.2 concerns the optional non-exact by-product; the author expressly permits skipping that section for the exact-clique proof.
